% Основные источники:
% https://sochisirius.ru/uploads/2021/11/math_1021_11a_12_gaussovy_celye_chisla_teoriya.pdf
% https://sochisirius.ru/uploads/2021/11/math_1021_11a_13_gaussovy_celye_chisla_zadachi.pdf
% https://sochisirius.ru/uploads/2021/11/math_1021_11b_19-celye_gaussovye_chisla.pdf
% https://sochisirius.ru/uploads/2021/11/math_1021_11v_14_gaussovy_chisla.pdf

\begin{center}
    {\LARGE\bfseries Целые гауссовы числа}
\end{center}

\begin{remarkbox}
Листок не предполагает предварительного знакомства с гауссовыми
числами. Можно пользоваться обычными свойствами целых чисел,
комплексным сопряжением и равенством $i^2=-1$.
\end{remarkbox}

\begin{defbox}{гауссовы целые числа и норма}
\emph{Целыми гауссовыми числами} называются комплексные числа вида
\[
    a+bi,\qquad a,b\in\mathbb Z.
\]
Множество всех таких чисел обозначается через
\[
    \mathbb Z[i]=\{a+bi\mid a,b\in\mathbb Z\}.
\]

Для $\alpha=a+bi$ положим
\[
    \overline{\alpha}=a-bi,
    \qquad
    N(\alpha)=\alpha\overline{\alpha}=a^2+b^2.
\]
Число $N(\alpha)$ называется \emph{нормой} числа $\alpha$.
\end{defbox}

\begin{center}
    {\large\bfseries I. Первые свойства}
\end{center}

\problem{1}{Докажите следующие свойства.}

\begin{enumerate}[label=\textbf{\alph*)}, leftmargin=8mm]
    \item Если $\alpha,\beta\in\mathbb Z[i]$, то
    \[
        \alpha+\beta,\quad \alpha-\beta,\quad
        \alpha\beta,\quad \overline{\alpha}
        \in\mathbb Z[i].
    \]

    \item Для любого $\alpha\in\mathbb Z[i]$ число $N(\alpha)$ является
    целым неотрицательным, причём
    \[
        N(\alpha)=0\quad\Longleftrightarrow\quad \alpha=0.
    \]

    \item Для любых $\alpha,\beta\in\mathbb Z[i]$ выполнено
    \[
        N(\alpha\beta)=N(\alpha)N(\beta).
    \]
\end{enumerate}

\medskip

\begin{defbox}{делимость, обратимые и ассоциированные элементы}
Для $\alpha,\beta\in\mathbb Z[i]$ говорят, что $\alpha$ делит $\beta$,
и пишут $\alpha\mid\beta$, если существует $\gamma\in\mathbb Z[i]$,
для которого
\[
    \beta=\alpha\gamma.
\]

Число $\varepsilon\in\mathbb Z[i]$ называется \emph{обратимым}, если
$\varepsilon^{-1}\in\mathbb Z[i]$.

Числа $\alpha$ и $\beta$ называются \emph{ассоциированными}, если
$\alpha=\varepsilon\beta$ для некоторого обратимого
$\varepsilon\in\mathbb Z[i]$.
\end{defbox}

\problem{2}{Исследуйте обратимые и ассоциированные элементы.}

\begin{enumerate}[label=\textbf{\alph*)}, leftmargin=8mm]
    \item Докажите, что $\varepsilon\in\mathbb Z[i]$ обратимо тогда
    и только тогда, когда
    \[
        N(\varepsilon)=1.
    \]

    \item Найдите все обратимые элементы кольца $\mathbb Z[i]$.

    \item Опишите геометрически все числа, ассоциированные с
    $\alpha=a+bi$.

    \item Верно ли, что два гауссовых числа одинаковой нормы обязательно
    ассоциированы?
\end{enumerate}

\medskip

\problem{3}{Исследуйте делимость на число $1+i$.}

\begin{enumerate}[label=\textbf{\alph*)}, leftmargin=8mm]
    \item Докажите, что
    \[
        1+i\mid a+bi
        \quad\Longleftrightarrow\quad
        a\equiv b\pmod 2.
    \]

    \item Докажите эквивалентность
    \[
        1+i\mid\alpha
        \quad\Longleftrightarrow\quad
        2\mid N(\alpha).
    \]

    \item Для каждого из чисел
    \[
        8+12i,\qquad 7+i,\qquad 5+3i
    \]
    найдите наибольшее $k$, для которого $(1+i)^k$ делит это число.
\end{enumerate}

\medskip

\problem{4}{Потренируйтесь работать с нормой и делителями.}

\begin{enumerate}[label=\textbf{\alph*)}, leftmargin=8mm]
    \item Найдите все гауссовы числа норм
    \[
        2,\qquad 5,\qquad 13,\qquad 17.
    \]

    \item Разложите числа $2$, $5$, $13$ и $17$ в произведение
    необратимых гауссовых чисел.

    \item Найдите с точностью до ассоциированности все делители числа
    $5$ в кольце $\mathbb Z[i]$.

    \item Докажите, что из $\alpha\mid\beta$, где $\alpha\ne0$, следует
    \[
        N(\alpha)\mid N(\beta).
    \]
    Покажите на примере чисел $1+2i$ и $2+i$, что обратное утверждение
    неверно.
\end{enumerate}

\begin{center}
    {\large\bfseries II. Деление с остатком и алгоритм Евклида}
\end{center}

\problem{5}{
Пусть $z=x+yi$ — произвольное комплексное число. Докажите, что
существует $q=m+ni\in\mathbb Z[i]$, для которого
\[
    |z-q|^2\leqslant\frac12.
\]
}

\textit{Указание.} Выберите $m$ и $n$ ближайшими целыми числами к
$x$ и $y$ соответственно.

\medskip

\problem{6}{
Докажите теорему о делении с остатком: для любых
$\alpha,\beta\in\mathbb Z[i]$, где $\beta\ne0$, существуют
$q,r\in\mathbb Z[i]$ такие, что
\[
    \alpha=\beta q+r,
    \qquad
    N(r)<N(\beta).
\]
}

\textit{Указание.} Примените предыдущую задачу к комплексному числу
$\alpha/\beta$.

\medskip

\problem{7}{Выполните деление с остатком.}

\begin{enumerate}[label=\textbf{\alph*)}, leftmargin=8mm]
    \item Разделите $17+9i$ на $4+3i$.

    \item Разделите $9+2i$ на $4+7i$.

    \item Найдите два различных допустимых остатка при делении
    $3+i$ на $1+i$.
\end{enumerate}

\begin{remarkbox}
В отличие от деления целых чисел, неполное частное и остаток
в $\mathbb Z[i]$ могут определяться неоднозначно.
\end{remarkbox}

\begin{defbox}{наибольший общий делитель}
Гауссово число $d$ называется \emph{наибольшим общим делителем}
чисел $\alpha$ и $\beta$, если
\[
    d\mid\alpha,\qquad d\mid\beta,
\]
и каждый общий делитель $\alpha$ и $\beta$ делит $d$.

Наибольший общий делитель определён с точностью до умножения на
обратимый элемент.
\end{defbox}

\problem{8}{Исследуйте алгоритм Евклида в $\mathbb Z[i]$.}

\begin{enumerate}[label=\textbf{\alph*)}, leftmargin=8mm]
    \item Докажите, что последовательное деление с остатком позволяет
    найти наибольший общий делитель двух гауссовых чисел.

    \item Найдите
    \[
        \gcd(4+7i,\;9+2i).
    \]

    \item Представьте найденный наибольший общий делитель в виде
    \[
        (4+7i)u+(9+2i)v,
        \qquad u,v\in\mathbb Z[i].
    \]

    \item Найдите
    \[
        \gcd(3+8i,\;7-5i).
    \]
\end{enumerate}

\medskip

\begin{defbox}{простое гауссово число}
Ненулевое необратимое число $\pi\in\mathbb Z[i]$ называется
\emph{простым гауссовым числом}, если из равенства
\[
    \pi=\alpha\beta
\]
следует, что хотя бы одно из чисел $\alpha,\beta$ обратимо.
\end{defbox}

\problem{9}{Докажите аналоги стандартных арифметических фактов.}

\begin{enumerate}[label=\textbf{\alph*)}, leftmargin=8mm]
    \item Наибольший общий делитель любых $\alpha,\beta\in\mathbb Z[i]$
    можно представить в виде
    \[
        d=\alpha u+\beta v,
        \qquad u,v\in\mathbb Z[i].
    \]

    \item Если $\gcd(\alpha,\beta)=1$ и
    \[
        \alpha\mid\beta\gamma,
    \]
    то $\alpha\mid\gamma$.

    \item Если $\pi$ — простое гауссово число и
    \[
        \pi\mid\alpha\beta,
    \]
    то $\pi\mid\alpha$ или $\pi\mid\beta$.

    \item Докажите существование разложения любого ненулевого
    необратимого гауссова числа в произведение простых.

    \item Докажите единственность такого разложения с точностью до
    порядка сомножителей и замены сомножителей ассоциированными.
\end{enumerate}

\begin{factbox}{основная теорема арифметики в $\mathbb Z[i]$}
Любое ненулевое гауссово число раскладывается в произведение простых
гауссовых чисел. Такое разложение единственно с точностью до порядка
сомножителей и умножения сомножителей на обратимые элементы.

Этим фактом можно пользоваться во всех последующих задачах.
\end{factbox}

\begin{center}
    {\large\bfseries III. Какие гауссовы числа являются простыми?}
\end{center}

\problem{10}{
Докажите, что если $N(\alpha)$ является обычным простым числом,
то $\alpha$ является простым гауссовым числом.
}

\medskip

\problem{11}{
Пусть $p$ — обычное простое число вида $4k+3$. Докажите, что
$p$ является простым числом и в кольце $\mathbb Z[i]$.
}

\textit{Указание.} Если $p=\alpha\beta$, рассмотрите возможные значения
$N(\alpha)$ и $N(\beta)$ и вспомните остатки квадратов по модулю $4$.

\medskip

\begin{remarkbox}
В следующей задаче можно пользоваться теоремами Ферма и Вильсона:
для простого $p$ выполнено
\[
    a^{p-1}\equiv1\pmod p\quad (p\nmid a),
    \qquad
    (p-1)!\equiv-1\pmod p.
\]
\end{remarkbox}

\problem{12}{
Пусть $p$ — обычное простое число вида $4k+1$.
}

\begin{enumerate}[label=\textbf{\alph*)}, leftmargin=8mm]
    \item Используя теорему Вильсона, докажите, что существует целое
    число $x$, для которого
    \[
        x^2\equiv-1\pmod p.
    \]

    \item В кольце $\mathbb Z[i]$ рассмотрите
    \[
        d=\gcd(p,\;x+i).
    \]
    Докажите, что $d$ не является ни обратимым, ни ассоциированным
    с $p$.

    \item Докажите, что
    \[
        N(d)=p.
    \]

    \item Выведите, что существуют натуральные числа $a,b$, для которых
    \[
        p=a^2+b^2,
        \qquad
        p=(a+bi)(a-bi).
    \]
\end{enumerate}

\medskip

\problem{13}{
Докажите, что с точностью до ассоциированности список всех простых
гауссовых чисел имеет следующий вид:
}

\begin{enumerate}[label=\textbf{\arabic*)}, leftmargin=10mm]
    \item число $1+i$;

    \item обычные простые числа $p\equiv3\pmod4$;

    \item числа $a+bi$, норма которых
    \[
        a^2+b^2
    \]
    является обычным простым числом вида $4k+1$.
\end{enumerate}

\begin{factbox}{разложение обычных простых чисел в $\mathbb Z[i]$}
С точностью до ассоциированности:
\[
\begin{aligned}
    2&=-i(1+i)^2,\\
    p&\text{ остаётся простым в }\mathbb Z[i],
        &&\text{если }p\equiv3\pmod4,\\
    p&=(a+bi)(a-bi),
        &&\text{если }p\equiv1\pmod4
          \text{ и }p=a^2+b^2.
\end{aligned}
\]
\end{factbox}

\begin{center}
    {\large\bfseries IV. Приложения}
\end{center}

\problem{14}{
Докажите, что если простое число $p\equiv1\pmod4$ представлено в виде
\[
    p=a^2+b^2,
\]
то это представление единственно с точностью до перестановки $a,b$
и изменения их знаков.
}

\medskip

\problem{15}{
Опишите все примитивные пифагоровы тройки. Иными словами, докажите,
что если
\[
    a^2+b^2=c^2,
    \qquad
    \gcd(a,b,c)=1,
\]
то после возможной перестановки $a$ и $b$ существуют взаимно простые
натуральные числа $m>n$ разной чётности, для которых
\[
    a=m^2-n^2,\qquad
    b=2mn,\qquad
    c=m^2+n^2.
\]
}

\textit{Указание.} Рассмотрите разложение
\[
    (a+bi)(a-bi)=c^2
\]
в кольце $\mathbb Z[i]$.

\medskip

\problem{16}{Докажите теорему о сумме двух квадратов.}

\begin{enumerate}[label=\textbf{\alph*)}, leftmargin=8mm]
    \item Получите из равенства
    \[
        (a+bi)(c+di)=(ac-bd)+(ad+bc)i
    \]
    тождество
    \[
        (a^2+b^2)(c^2+d^2)
        =(ac-bd)^2+(ad+bc)^2.
    \]

    \item Докажите, что натуральное число $n$ представимо в виде
    \[
        n=x^2+y^2,\qquad x,y\in\mathbb Z,
    \]
    тогда и только тогда, когда каждый простой делитель
    $p\equiv3\pmod4$ входит в разложение числа $n$ в чётной степени.
\end{enumerate}

\medskip

\problem{17}{
Решите в целых числах уравнение
\[
    x^2+4=y^3.
\]
}

\textit{Указание.} После разбора чётности рассмотрите в
$\mathbb Z[i]$ равенство
\[
    (x+2i)(x-2i)=y^3
\]
и исследуйте наибольший общий делитель множителей.

% Источник: листок 11А «Гауссовы целые числа. Задачи», задача 3а.

\medskip

\problem{18}{
Натуральные числа $x,y,z$ удовлетворяют равенству
\[
    xy=z^2+1.
\]
Докажите, что существуют целые числа $a,b,c,d$, для которых
\[
    x=a^2+b^2,\qquad
    y=c^2+d^2,\qquad
    z=ac+bd.
\]
}

\textit{Указание.} Разложите число $z+i$ на простые множители в
$\mathbb Z[i]$ и распределите их между двумя множителями с нормами
$x$ и $y$.

% Источник: листок 11А «Гауссовы целые числа. Задачи», задача 4.

\medskip

\problem{19}{
Найдите все пары простых чисел $p$ и $q$ (не обязательно различных),
для которых каждое из чисел
\[
    2p-1,\qquad 2q-1,\qquad 2pq-1
\]
является квадратом натурального числа.
}

% Источник: листки Сириуса для групп 11А, 11Б и 11В.